\section{The fold: existence, multiplicity and the sizing iteration}
\label{sec:fold}

This section analyses \cref{eq:closure-q} completely. The results are
elementary but they are the substance of the problem, and they are what
separates ``the design is heavy'' from ``the design does not exist''.

Throughout, fix $\alpha>0$, $\beta>0$, $\Mzero>0$ and $q>0$, and define
\begin{equation}
  F:(0,\infty)\to\R, \qquad F(m) \coloneqq \alpha m - \beta m^{q} .
  \label{eq:F}
\end{equation}
$F(m)$ is the fixed load that a machine of gross mass $m$ can support: gross
mass less what structure, propulsion and battery consume. Solving the closure
means solving $F(m)=\Mzero$.

\subsection{The trichotomy}

\begin{lemma}[Strict concavity]
\label{lem:concave}
If $q>1$ then $F$ is strictly concave on $(0,\infty)$, attains a unique
maximum at
\begin{equation}
  m^{*} = \left(\frac{\alpha}{q\beta}\right)^{\frac{1}{q-1}} ,
  \label{eq:mstar}
\end{equation}
and
\begin{equation}
  \Mmax \coloneqq F(m^{*}) = \alpha\, m^{*}\left(1 - \frac{1}{q}\right) .
  \label{eq:Mmax}
\end{equation}
Moreover $F(0^+)=0$, $F$ is strictly increasing on $(0,m^*)$, strictly
decreasing on $(m^*,\infty)$, and $F(m)\to-\infty$ as $m\to\infty$.
\end{lemma}

\begin{proof}
$F$ is smooth on $(0,\infty)$ with
\begin{equation}
  F'(m) = \alpha - q\beta m^{q-1},
  \qquad
  F''(m) = -q(q-1)\beta m^{q-2} .
  \label{eq:Fderiv}
\end{equation}
For $q>1$ and $m>0$ we have $q(q-1)\beta m^{q-2}>0$, hence $F''<0$ everywhere
on $(0,\infty)$: $F$ is strictly concave. Setting $F'(m)=0$ gives
$m^{q-1} = \alpha/(q\beta)$, and since $q-1>0$ the map $m\mapsto m^{q-1}$ is a
strictly increasing bijection of $(0,\infty)$ onto itself, so the stationary
point is unique and given by \cref{eq:mstar}. Strict concavity makes it the
global maximum. Evaluating,
\[
  F(m^*) = \alpha m^* - \beta (m^*)^{q}
         = \alpha m^* - \beta m^* (m^*)^{q-1}
         = \alpha m^* - \beta m^* \frac{\alpha}{q\beta}
         = \alpha m^*\left(1-\frac1q\right),
\]
which is \cref{eq:Mmax}. The sign of $F'$ follows from monotonicity of
$m^{q-1}$, and $F(m) = m(\alpha - \beta m^{q-1}) \to -\infty$ because the
bracket tends to $-\infty$.
\end{proof}

\begin{theorem}[Trichotomy for the closure]
\label{thm:trichotomy}
Let $q>1$, $\alpha,\beta,\Mzero>0$, and let $\Mmax$ be as in \cref{eq:Mmax}.
Then the equation $F(m)=\Mzero$ has, on $(0,\infty)$,
\begin{enumerate}[label=(\roman*),nosep]
  \item exactly two solutions $0 < m_- < m^* < m_+$ if $\Mzero < \Mmax$;
  \item exactly one solution $m = m^*$, a double root, if $\Mzero = \Mmax$;
  \item no solution if $\Mzero > \Mmax$.
\end{enumerate}
\end{theorem}

\begin{proof}
By \cref{lem:concave}, $F$ increases strictly from $F(0^+)=0$ to $F(m^*)=\Mmax$
on $(0,m^*]$ and decreases strictly from $\Mmax$ to $-\infty$ on
$[m^*,\infty)$. Both branches are continuous and strictly monotone.

(iii) If $\Mzero>\Mmax$ then $F(m)\le \Mmax < \Mzero$ for all $m>0$, so no
solution exists.

(ii) If $\Mzero=\Mmax$, then $F(m)=\Mzero$ forces $m=m^*$ by uniqueness of the
maximiser. It is a double root because $F'(m^*)=0$, so $F(m)-\Mzero$ has a zero
of order at least two; the order is exactly two since $F''(m^*)<0$.

(i) If $0<\Mzero<\Mmax$ then on $(0,m^*]$ the function $F$ is a strictly
increasing continuous bijection onto $(0,\Mmax]$, so there is exactly one
$m_-\in(0,m^*)$ with $F(m_-)=\Mzero$. On $[m^*,\infty)$, $F$ is a strictly
decreasing continuous bijection onto $(-\infty,\Mmax]$, so there is exactly one
$m_+\in(m^*,\infty)$ with $F(m_+)=\Mzero$. Strictness gives $m_-<m^*<m_+$ and
excludes further roots.
\end{proof}

\begin{corollary}[Closed forms at $q=3/2$]
\label{cor:closedform}
For the fixed-rotor case $q=3/2$,
\begin{equation}
  m^{*} = \frac{4\alpha^{2}}{9\beta^{2}},
  \qquad
  \Mmax = \frac{4\alpha^{3}}{27\beta^{2}} .
  \label{eq:closedform}
\end{equation}
\end{corollary}

\begin{proof}
Put $q=3/2$ in \cref{eq:mstar}: $1/(q-1)=2$ and $\alpha/(q\beta)
= 2\alpha/(3\beta)$, so $m^* = 4\alpha^2/(9\beta^2)$. Then \cref{eq:Mmax} gives
$\Mmax = \alpha m^*(1-2/3) = \alpha m^*/3 = 4\alpha^3/(27\beta^2)$.
\end{proof}

\begin{remark}[Which root is the design]
The physically meaningful root is $m_-$. Both roots satisfy the mass budget,
but $m_+$ is a machine that is heavier than necessary, carries a battery sized
for its own excess weight, and is, by \cref{thm:iteration} below, unstable
under the natural sizing iteration. It is a mathematical artefact of the
quadratic-like structure, not a design.
\end{remark}

\begin{remark}[Cubic reduction]
\Cref{eq:closure} at $q=3/2$ becomes a cubic under the substitution
$u=\sqrt m$: $\beta u^3 - \alpha u^2 + \Mzero = 0$. The product of its three
roots is $-\Mzero/\beta<0$, so exactly one root is negative and the other two
are the pair $\sqrt{m_\pm}$ of \cref{thm:trichotomy}, real precisely when the
discriminant is non-negative. The engine solves the equation by bracketing
rather than by radicals, for reasons given in \cref{sec:numerics}.
\end{remark}

\subsection{The transition is a saddle-node bifurcation}

Regard $\Mzero$ as a parameter and consider the one-parameter family of
equations $G(m;\Mzero) \coloneqq F(m)-\Mzero = 0$.

\begin{theorem}[Saddle-node]
\label{thm:bifurcation}
At $(m,\Mzero)=(m^*,\Mmax)$ the family $G$ undergoes a saddle-node
(fold) bifurcation: the defining conditions
\begin{equation}
  G = 0, \qquad \frac{\partial G}{\partial m} = 0
  \label{eq:sn1}
\end{equation}
hold, together with the nondegeneracy and transversality conditions
\begin{equation}
  \frac{\partial^2 G}{\partial m^2}(m^*,\Mmax) = -q(q-1)\beta (m^*)^{q-2} \ne 0,
  \qquad
  \frac{\partial G}{\partial \Mzero}(m^*,\Mmax) = -1 \ne 0 .
  \label{eq:sn2}
\end{equation}
Consequently, near the bifurcation the two roots behave as
\begin{equation}
  m_{\pm} = m^* \pm \sqrt{\frac{2(\Mmax-\Mzero)}{q(q-1)\beta (m^*)^{q-2}}}
            + O(\Mmax-\Mzero) ,
  \label{eq:sqrtlaw}
\end{equation}
so they approach one another with a square-root law and annihilate.
\end{theorem}

\begin{proof}
\Cref{eq:sn1} is \cref{lem:concave} together with the definition of $\Mmax$.
\Cref{eq:sn2} follows from \cref{eq:Fderiv} and $\partial G/\partial \Mzero=-1$.
These are exactly the hypotheses of the saddle-node normal form theorem
(see e.g.\ Guckenheimer and Holmes~\cite[\S3.4]{guckenheimer1983} or
Kuznetsov~\cite[\S3.2]{kuznetsov2004}), which asserts the existence of a smooth
change of coordinates taking $G$ to $\dot y = \mu - y^2$ near the point.
For \cref{eq:sqrtlaw}, Taylor-expand $F$ about $m^*$:
\[
  \Mzero = F(m) = \Mmax + \tfrac12 F''(m^*)(m-m^*)^2 + O((m-m^*)^3),
\]
using $F'(m^*)=0$. Solving for $m-m^*$ and inserting
$F''(m^*)=-q(q-1)\beta(m^*)^{q-2}$ gives the stated expansion.
\end{proof}

\begin{remark}[Why this matters practically]
\Cref{eq:sqrtlaw} says that near the boundary of feasibility the design mass is
an infinitely sensitive function of the requirement: $\dd m_-/\dd \Mzero \to
\infty$ as $\Mzero\uparrow\Mmax$. A design placed close to the fold is not
merely marginal; its mass is unbounded in its derivative with respect to every
assumption. This is a stronger statement than ``leave margin'' and is the
reason the engine reports the utilisation $\Mzero/\Mmax$ as a first-class
quantity.
\end{remark}

\subsection{The naive sizing iteration}

The obvious way to size such a vehicle is to guess a mass, compute what it
implies, and repeat:
\begin{equation}
  m_{k+1} = \Phi(m_k),
  \qquad
  \Phi(m) \coloneqq \Mzero + f_s m + \gamma_m m + \beta m^{q} .
  \label{eq:iteration}
\end{equation}
$\Phi$ is exactly \cref{eq:budget} read as an assignment. Its fixed points are
the roots of \cref{eq:closure-q}, since
$\Phi(m)=m \iff F(m)=\Mzero$.

\begin{theorem}[Convergence of the sizing iteration]
\label{thm:iteration}
Let $q>1$ and $\Mzero<\Mmax$, with roots $m_-<m^*<m_+$ as in
\cref{thm:trichotomy}. Then
\begin{enumerate}[label=(\roman*),nosep]
  \item $\Phi'(m) = 1 - F'(m)$, so $\Phi'(m_-)<1$ and $\Phi'(m_+)>1$;
  \item $m_-$ is an attracting fixed point of \cref{eq:iteration} with
        asymptotic linear rate $\Phi'(m_-) = 1-F'(m_-) \in (0,1)$, and $m_+$ is
        repelling;
  \item the iteration started from any $m_0 \in [0,m_+)$ converges
        monotonically upward to $m_-$ if $m_0<m_-$, and monotonically downward
        to $m_-$ if $m_0\in(m_-,m_+)$;
  \item for $m_0>m_+$ the iterates increase without bound;
  \item if $\Mzero>\Mmax$ the iterates increase without bound from every
        starting point.
\end{enumerate}
\end{theorem}

\begin{proof}
(i) Differentiating \cref{eq:iteration},
$\Phi'(m) = f_s+\gamma_m + q\beta m^{q-1}$. Using
$f_s+\gamma_m = 1-\alpha$ from \cref{eq:alpha} and
$F'(m)=\alpha - q\beta m^{q-1}$ from \cref{eq:Fderiv},
\[
  \Phi'(m) = 1-\alpha + q\beta m^{q-1} = 1 - F'(m).
\]
By \cref{lem:concave}, $F'(m_-)>0$ and $F'(m_+)<0$, giving $\Phi'(m_-)<1$ and
$\Phi'(m_+)>1$. Also $\Phi'>0$ everywhere since $f_s,\gamma_m,\beta\ge0$ and not
all zero, so $\Phi$ is strictly increasing; and $\Phi'(m_-) = 1-F'(m_-) > 0$
because $F'(m_-) \le \alpha < 1$.

(ii) Standard: a fixed point of a $C^1$ map is attracting when
$|\Phi'|<1$ and repelling when $|\Phi'|>1$, with linear rate $|\Phi'|$.

(iii) $\Phi$ is increasing and $\Phi(m)-m = \Mzero-F(m)$, which is negative
exactly on $(m_-,m_+)$ and positive on $[0,m_-)$. Hence for $m_0<m_-$ the
sequence is increasing and bounded above by $m_-$ (because $\Phi$ increasing
and $\Phi(m_-)=m_-$ imply $m_k<m_-\Rightarrow m_{k+1}<m_-$), so it converges,
necessarily to a fixed point in $[m_0,m_-]$, which must be $m_-$. The argument
on $(m_-,m_+)$ is symmetric with the sequence decreasing.

(iv) For $m_0>m_+$, $\Phi(m)-m = \Mzero - F(m) > 0$ since $F(m)<\Mzero$ there,
so the sequence increases; it cannot converge because there is no fixed point
above $m_+$; hence it diverges.

(v) If $\Mzero>\Mmax$ then $\Mzero-F(m)>0$ for every $m$, so
$m_{k+1}>m_k$ always and there is no fixed point to converge to.
\end{proof}

\begin{corollary}[Local stability criterion]
\label{cor:dmdm}
A design is a stable fixed point of the sizing iteration if and only if
\begin{equation}
  \frac{\dd}{\dd m}\Big[\mstr(m)+\mprop(m)+\mbat(m)\Big] < 1 ,
  \label{eq:dmdm}
\end{equation}
which for the model of \cref{sec:closure} reads
$f_s + \gamma_m + q\beta m^{q-1} < 1$, equivalently $F'(m)>0$, equivalently
$m<m^*$.
\end{corollary}

\Cref{cor:dmdm} is the precise form of the intuition that ``each extra kilogram
must require less than a kilogram of extra machine to carry it''. The engine
evaluates \cref{eq:dmdm} numerically at the solved design point and reports it.

\begin{remark}[Critical slowing down]
By \cref{thm:iteration}(ii) the convergence rate is $1-F'(m_-)$, and
$F'(m_-)\to 0$ as $\Mzero\uparrow\Mmax$. The iteration therefore becomes
arbitrarily slow near the fold: near-infeasible designs are not distinguished
from infeasible ones by a fixed iteration budget. This is a real trap. A
practical consequence is that a solver which reports ``did not converge in $K$
iterations'' cannot be used to decide feasibility; the analytic criterion
$\Mzero \lessgtr \Mmax$ must be used instead. The engine does exactly this and
classifies slow convergence separately from divergence.
\end{remark}

\subsection{Behaviour when \texorpdfstring{$q\le1$}{q <= 1}}

\begin{theorem}[No fold for $q\le1$]
\label{thm:pfold}
Let $\alpha,\beta,\Mzero>0$.
\begin{enumerate}[label=(\roman*),nosep]
  \item If $q=1$ and $\alpha>\beta$, then $F(m)=\Mzero$ has the unique solution
        $m=\Mzero/(\alpha-\beta)$; if $\alpha\le\beta$ there is no solution and
        $F\le 0$.
  \item If $0<q<1$, then $F(m)=\Mzero$ has exactly one solution for every
        $\Mzero>0$. In particular no critical payload exists.
\end{enumerate}
Consequently, under the disk-area scaling \cref{eq:areascaling} a fold exists
if and only if $p<1$.
\end{theorem}

\begin{proof}
(i) $F(m)=(\alpha-\beta)m$ is linear; the claim is immediate.

(ii) For $0<q<1$, $F'(m)=\alpha-q\beta m^{q-1}$ with $q-1<0$, so
$m^{q-1}\to\infty$ as $m\to0^+$ and $m^{q-1}\to0$ as $m\to\infty$. Hence
$F'(m)<0$ for small $m$ and $F'(m)>0$ for large $m$, with a single sign change
at $m_{\min}=(q\beta/\alpha)^{1/(1-q)}$. Thus $F$ decreases from $F(0^+)=0$ to
a negative minimum and then increases strictly to $+\infty$ (since
$F(m)=m(\alpha-\beta m^{q-1})$ and the bracket tends to $\alpha>0$). On
$[m_{\min},\infty)$, $F$ is a strictly increasing continuous bijection onto
$[F(m_{\min}),\infty)\supset(0,\infty)$, so there is exactly one root with
$F=\Mzero>0$, and none on $(0,m_{\min})$ where $F\le0<\Mzero$.

The final statement follows from $q=(3-p)/2$: $q>1 \iff p<1$.
\end{proof}

\begin{finding}
The catastrophe is not a property of flight. It is a property of holding the
rotor area fixed while the mass grows. Let the disk grow at least as fast as
the mass and the fold disappears entirely, replaced by the far gentler
condition \cref{eq:constdlfeas}. What is then unbounded is not the mass but the
rotor diameter, so the binding constraint migrates from energy to geometry,
noise and safety. It does not go away. Nor does the absence of a fold mean the
absence of a ceiling: at $p=1$ exactly the longest mission that closes is still
finite (\cref{prop:twall}); only for $p>1$ does it disappear.
\end{finding}
